\documentclass{article}
\usepackage[utf8]{inputenc}
\usepackage{amsmath}
\usepackage{algorithm}
\usepackage{algorithmic}
\usepackage{hyperref}

\title{Human-Model Alignment Analysis Pipeline}
\author{}
\date{}

\begin{document}

\maketitle

\section{Overview}

This document describes the comprehensive pipeline for analyzing human-model alignment in vision-language tasks. The pipeline processes human responses and model outputs, computes evaluation metrics, and generates comparative visualizations to assess alignment quality.

\section{Pipeline Architecture}

\subsection{Data Sources}

\begin{itemize}
    \item \textbf{Human Responses}: Collected from crowdsourcing experiments with confidence ratings
    \begin{itemize}
        \item VQA (open-ended text): $n=10$ responses per question from VQA v2 validation set
        \item Multiple Choice: $n=10$ responses per question from MMStar dataset
    \end{itemize}
    \item \textbf{Model Outputs}: Inference results from vision-language models
    \begin{itemize}
        \item InternVL3.5 family (1B, 2B, 4B, 8B parameters)
        \item Qwen3-VL family (2B, 4B, 8B parameters)
        \item Llava family (7B parameters)
    \end{itemize}
    \item \textbf{Ground Truth}: VQA v2 annotations (10 human annotators per question)
\end{itemize}

\subsection{Processing Stages}

\subsubsection{Stage 1: Human Response Processing}

\textbf{Script}: \texttt{score\_human\_results.py}

\textbf{Input}: Training data files containing aggregated human responses with confidence distributions

\textbf{Process}:
\begin{enumerate}
    \item Load human responses from JSONL files
    \item Extract consensus answers (highest confidence)
    \item For VQA (text answers):
    \begin{itemize}
        \item Normalize answers: lowercase, remove articles, strip punctuation
        \item Compute VQA accuracy: $\text{acc} = \min(1, \frac{\text{\# matches}}{3})$
        \item Compute embedding similarity (optional): cosine similarity using SentenceTransformer
    \end{itemize}
    \item For MC (choice answers):
    \begin{itemize}
        \item Extract letter choice (A/B/C/D)
        \item Compare with ground truth
    \end{itemize}
    \item Aggregate statistics:
    \begin{itemize}
        \item Confidence distributions
        \item Answer frequency distributions
        \item Question type distributions
        \item Mean/std of accuracy and similarity scores
    \end{itemize}
    \item Save QID mappings for each answer type
\end{enumerate}

\textbf{Output}:
\begin{itemize}
    \item \texttt{human\_vqa\_scored.jsonl}: Scored VQA responses
    \item \texttt{human\_mc\_scored.jsonl}: Scored MC responses
    \item \texttt{human\_vqa\_qids.json}: Question ID mappings for VQA
    \item \texttt{human\_mc\_qids.json}: Question ID mappings for MC
    \item \texttt{human\_vqa\_stats.json}: Aggregate statistics for VQA
    \item \texttt{human\_mc\_stats.json}: Aggregate statistics for MC
\end{itemize}

\subsubsection{Stage 2: Model Output Processing}

\textbf{Script}: \texttt{score\_results.py}

\textbf{Input}: Model inference JSONL files from multiple models and datasets

\textbf{Process}:
\begin{enumerate}
    \item Parse filename to extract model, dataset, and condition
    \item For each model output:
    \begin{itemize}
        \item Remove thinking tags: \texttt{<think>...</think>}
        \item For VQA: Normalize and compute VQA accuracy
        \item For MC: Extract choice letter using regex patterns
        \item (Optional) Compute embedding similarity for VQA
    \end{itemize}
    \item Compute per-category accuracy breakdowns
    \item Save processed results with scores
\end{enumerate}

\textbf{Output}:
\begin{itemize}
    \item Scored JSONL files in \texttt{evaluation/data/scored/}
    \item Summary CSV with model-wise performance
    \item Aggregate results JSON with category breakdowns
\end{itemize}

\subsubsection{Stage 3: Comparative Analysis \& Visualization}

\textbf{Script}: \texttt{visualize\_human\_analysis.py}

\textbf{Input}: Scored human and model results

\textbf{Visualizations Generated}:

\begin{enumerate}
    \item \textbf{Confidence Distribution}
    \begin{itemize}
        \item Bar plot of human confidence ratings (mapped from 1-5 scale to 0.05-1.0)
        \item Shows distribution of certainty across responses
    \end{itemize}

    \item \textbf{Accuracy Comparison}
    \begin{itemize}
        \item Histogram: Density plots of accuracy distributions (human vs models)
        \item Box plot: Statistical comparison of accuracy spread
        \item Reveals alignment in correctness patterns
    \end{itemize}

    \item \textbf{Answer Similarity Comparison}
    \begin{itemize}
        \item Histogram: Embedding similarity distributions
        \item CDF plot: Cumulative distribution functions for human and models
        \item Measures semantic alignment beyond binary correctness
    \end{itemize}

    \item \textbf{Question Type Distribution}
    \begin{itemize}
        \item Bar plot by question word (what, where, when, who, how, etc.)
        \item Identifies which question types humans were asked
    \end{itemize}

    \item \textbf{Confidence vs Accuracy Relationship}
    \begin{itemize}
        \item Scatter plot with trend line
        \item Binned averages showing calibration quality
        \item Tests whether higher confidence correlates with higher accuracy
    \end{itemize}

    \item \textbf{Answer Frequency Distribution}
    \begin{itemize}
        \item Top 20 most common answers
        \item Reveals response patterns and biases
    \end{itemize}
\end{enumerate}

\textbf{Output}: High-resolution PNG figures in \texttt{evaluation/figures/}

\section{Evaluation Metrics}

\subsection{VQA Accuracy}

Standard VQA metric following Das et al. (2017):

\begin{equation}
\text{Acc}_{\text{VQA}} = \min\left(1, \frac{\#\{\text{humans that said answer}\}}{3}\right)
\end{equation}

An answer is considered correct if at least 3 out of 10 ground truth annotators provided the same answer.

\subsection{Multiple Choice Accuracy}

Binary accuracy for extracted choices:

\begin{equation}
\text{Acc}_{\text{MC}} = \begin{cases}
1 & \text{if extracted choice} = \text{ground truth} \\
0 & \text{otherwise}
\end{cases}
\end{equation}

\subsection{Answer Similarity}

Semantic similarity using SentenceTransformer (all-MiniLM-L6-v2):

\begin{equation}
\text{Sim}(a_{\text{pred}}, a_{\text{gt}}) = \frac{\mathbf{e}_{\text{pred}} \cdot \mathbf{e}_{\text{gt}}}{\|\mathbf{e}_{\text{pred}}\| \|\mathbf{e}_{\text{gt}}\|}
\end{equation}

where $\mathbf{e}$ are embedding vectors. For VQA with multiple ground truth answers, we use the majority answer.

\subsection{Confidence Mapping}

Human confidence ratings are mapped to continuous values:

\begin{table}[h]
\centering
\begin{tabular}{|c|c|}
\hline
\textbf{Rating} & \textbf{Confidence Value} \\
\hline
1 (Very Uncertain) & 0.05 \\
2 (Uncertain) & 0.25 \\
3 (Neutral) & 0.50 \\
4 (Confident) & 0.75 \\
5 (Very Confident) & 1.00 \\
\hline
yes (Quality Check) & 1.00 \\
maybe & 0.50 \\
no & 0.01 \\
\hline
\end{tabular}
\caption{Confidence rating mappings}
\end{table}

\section{Usage Instructions}

\subsection{Complete Pipeline Execution}

\begin{algorithm}
\caption{Human-Model Alignment Analysis}
\begin{algorithmic}[1]
\STATE \textbf{Step 1}: Score human responses
\STATE \texttt{python score\_human\_results.py \\
    --text\_data data/training/s1\_text/train\_agg\_10\_blind\_inst.jsonl \\
    --choice\_data data/training/s1\_choice/train\_agg\_10\_blind\_inst.jsonl \\
    --output\_dir evaluation/human\_scored/ \\
    --with\_similarity}

\STATE \textbf{Step 2}: Score model outputs (if not already done)
\STATE \texttt{python score\_results.py \\
    --input\_dir data/models/ \\
    --output\_dir evaluation/data/scored/ \\
    --with\_similarity}

\STATE \textbf{Step 3}: Generate visualizations
\STATE \texttt{python visualize\_human\_analysis.py \\
    --human\_dir evaluation/human\_scored/ \\
    --model\_dir evaluation/data/scored/ \\
    --models InternVL3\_5-2B Qwen3-VL-4B-Instruct \\
    --dataset vqa\_1k\_inst\_blind \\
    --output\_dir evaluation/figures/ \\
    --answer\_type vqa}
\end{algorithmic}
\end{algorithm}

\subsection{Quick VQA-Only Analysis}

For VQA analysis without model comparison:

\begin{verbatim}
# Score human VQA responses
python score_human_results.py \
    --text_data data/training/s1_text/train_agg_10_blind_inst.jsonl \
    --output_dir evaluation/human_scored/ \
    --with_similarity

# Generate VQA visualizations
python visualize_human_analysis.py \
    --human_dir evaluation/human_scored/ \
    --output_dir evaluation/figures/ \
    --answer_type vqa
\end{verbatim}

\section{Expected Outputs}

\subsection{Scored Results Structure}

\textbf{Human VQA Result} (\texttt{human\_vqa\_scored.jsonl}):
\begin{verbatim}
{
  "qid": "418297001",
  "answer": "yes",
  "confidence": 0.75,
  "accuracy": 1.0,
  "correct": true,
  "gt_answers": ["yes", "yes", "yes", ...],
  "all_human_answers": ["yes", "no"],
  "all_confidences": [0.75, 0.25],
  "answer_similarity": 0.98,
  "question_type": "is"
}
\end{verbatim}

\textbf{Statistics File} (\texttt{human\_vqa\_stats.json}):
\begin{verbatim}
{
  "total_questions": 374,
  "total_responses": 3740,
  "mean_accuracy": 0.6234,
  "std_accuracy": 0.3421,
  "correct_count": 267,
  "mean_similarity": 0.7845,
  "std_similarity": 0.1923,
  "confidence_dist": {"0.05": 45, "0.25": 123, ...},
  "question_type_dist": {"what": 156, "is": 89, ...}
}
\end{verbatim}

\subsection{QID Mappings}

Question IDs are saved for later filtering and comparison:

\begin{verbatim}
{
  "qids": ["418297001", "444302001", ...],
  "count": 374
}
\end{verbatim}

These QIDs enable:
\begin{itemize}
    \item Filtering model results to match human-evaluated questions
    \item Fair comparison on identical question sets
    \item Tracking which questions were used for each task type
\end{itemize}

\section{Statistical Analyses}

\subsection{Distribution Comparisons}

For each metric (accuracy, similarity), we compute:

\begin{enumerate}
    \item \textbf{Descriptive Statistics}
    \begin{itemize}
        \item Mean ($\mu$), Standard Deviation ($\sigma$)
        \item Median, Quartiles (Q1, Q3)
        \item Min, Max, Range
    \end{itemize}

    \item \textbf{Distribution Tests}
    \begin{itemize}
        \item Kolmogorov-Smirnov test for distribution similarity
        \item Mann-Whitney U test for median differences
        \item Effect size (Cohen's d) for practical significance
    \end{itemize}

    \item \textbf{Correlation Analysis}
    \begin{itemize}
        \item Pearson correlation: confidence vs accuracy
        \item Spearman correlation: robustness to outliers
    \end{itemize}
\end{enumerate}

\subsection{Alignment Metrics}

\textbf{Response Overlap}: Percentage of questions where human consensus matches model output

\textbf{Confidence Calibration}: Expected Calibration Error (ECE) comparing confidence bins with accuracy

\textbf{Embedding Distance}: Average L2 distance in semantic space

\section{Key Findings Template}

When writing the paper, report:

\begin{enumerate}
    \item \textbf{Human Performance Baseline}
    \begin{itemize}
        \item Mean accuracy: $\mu_{\text{acc}} = X.XX \pm X.XX$
        \item Confidence-accuracy correlation: $r = X.XX$ ($p < 0.001$)
        \item Most common question types: [list top 3]
    \end{itemize}

    \item \textbf{Human-Model Alignment}
    \begin{itemize}
        \item Best-aligned model: [name] with similarity $X.XX$
        \item Accuracy gap: $\Delta_{\text{acc}} = |\mu_{\text{human}} - \mu_{\text{model}}|$
        \item Distribution overlap: KS statistic $= X.XX$ ($p = X.XX$)
    \end{itemize}

    \item \textbf{Confidence Analysis}
    \begin{itemize}
        \item Calibration quality: ECE $= X.XX$
        \item Over/under-confident patterns
        \item Confidence ranges by question type
    \end{itemize}
\end{enumerate}

\section{Troubleshooting}

\subsection{Common Issues}

\textbf{Missing VQA annotations}: Ensure \texttt{v2\_mscoco\_val2014\_annotations.json} is in correct path

\textbf{CUDA out of memory}: Use CPU encoder by modifying \texttt{get\_encoder()} to remove \texttt{.to('cuda')}

\textbf{Empty visualizations}: Check that answer\_type matches data (vqa vs mc)

\textbf{QID mismatches}: Verify question IDs are consistent across datasets (integer vs string format)

\section{Future Extensions}

\begin{itemize}
    \item \textbf{Category-wise Analysis}: Break down metrics by VQA question type (yes/no, number, color, etc.)
    \item \textbf{Difficulty Stratification}: Analyze alignment on easy vs hard questions
    \item \textbf{Multi-turn Analysis}: Extend to conversational settings
    \item \textbf{Uncertainty Quantification}: Model predictive uncertainty comparison with human confidence
\end{itemize}

\section{References}

\begin{itemize}
    \item VQA v2 Dataset: Goyal et al. (2017)
    \item VQA Evaluation Metric: Das et al. (2017)
    \item MMStar Dataset: Chen et al. (2024)
    \item SentenceTransformers: Reimers \& Gurevych (2019)
\end{itemize}

\end{document}
